% GENERATED FILE. Edit canonical lessons, then run npm run generate:books.
% canonical-source-hash: fnv1a64:ce7a030afb51efba
% canonical-lessons: UR-C269-jana-chahiye, UR-C269-chunanche, UR-C269-is-liye, UR-R269-advice-and-reasons

\chapter{Should, Must, and So}
\label{ch:ur-should-must-and-so}
\hlchaptermodality{269}

\begin{chapteropening}
\textbf{By the end of this chapter:} \emph{I can say what someone should do (āp ko jānā chāhiye), what I have to do (mujhe jānā paṛe gā), and what follows from a reason (is liye, to, chunāṅche).}

Complete UR-R269-advice-and-reasons: give advice, say what you have to do, and say what follows.
\end{chapteropening}

\section[\texorpdfstring{\ur{جانا} \ur{چاہیے} (jānā chāhiye)}{جانا چاہیے (jānā chāhiye)}]{\ur{جانا} \ur{چاہیے} (jānā chāhiye) — should, and have to}

\label{lesson:UR-C269-jana-chahiye}



\begin{quote}
Say \textquotedblleft{}I need water\textquotedblright{}. (\emph{mujhe pānī chāhiye}.) \textbf{\ur{چاہیے}} (\emph{chāhiye}) after a noun is \textquotedblleft{}is needed\textquotedblright{}. After a verb it is advice.
\end{quote}

\subsection*{You'll want to know: should}
\begin{quote}
\textbf{\ur{جانا} \ur{چاہیے۔}}
\end{quote}

\begin{quote}
\emph{jānā chāhiye.} — \textquotedblleft{}One should go.\textquotedblright{}
\end{quote}

\begin{quote}
\textbf{\ur{آپ} \ur{کو} \ur{جانا} \ur{چاہیے۔}}
\end{quote}

\begin{quote}
\emph{āp ko jānā chāhiye.} — \textquotedblleft{}You should go.\textquotedblright{}
\end{quote}

The one who should takes \textbf{\ur{کو}} (\emph{ko}), and the verb stays in its \textbf{-\ur{نا}} form:

\begin{quote}
\textbf{\ur{آپ} \ur{کو} \ur{آرام} \ur{کرنا} \ur{چاہیے۔}}
\end{quote}

\begin{quote}
\emph{āp ko ārām karnā chāhiye.} — \textquotedblleft{}You should rest.\textquotedblright{}
\end{quote}

\begin{quote}
\textbf{\ur{مجھے} \ur{پانی} \ur{پینا} \ur{چاہیے۔}}
\end{quote}

\begin{quote}
\emph{mujhe pānī pīnā chāhiye.} — \textquotedblleft{}I should drink water.\textquotedblright{}
\end{quote}

\begin{quote}
\textbf{\ur{آپ} \ur{کو} \ur{نہیں} \ur{جانا} \ur{چاہیے۔}}
\end{quote}

\begin{quote}
\emph{āp ko nahīṅ jānā chāhiye.} — \textquotedblleft{}You shouldn't go.\textquotedblright{}
\end{quote}

\subsection*{You'll want to know: have to}
You know \textbf{\ur{پڑنا}} (\emph{paṛnā}) from \textbf{\ur{بیمار} \ur{پڑنا}} (\emph{bīmār paṛnā}), \textquotedblleft{}to fall ill\textquotedblright{}. After a verb's \textbf{-\ur{نا}} form it is \textquotedblleft{}to have to\textquotedblright{}, and it carries the time:

\begin{quote}
\textbf{\ur{مجھے} \ur{جانا} \ur{پڑے} \ur{گا۔}}
\end{quote}

\begin{quote}
\emph{mujhe jānā paṛe gā.} — \textquotedblleft{}I'll have to go.\textquotedblright{}
\end{quote}

\begin{quote}
\textbf{\ur{مجھے} \ur{روز} \ur{کام} \ur{کرنا} \ur{پڑتا} \ur{ہے۔}}
\end{quote}

\begin{quote}
\emph{mujhe roz kām karnā paṛtā hai.} — \textquotedblleft{}I have to work every day.\textquotedblright{}
\end{quote}

\begin{quote}
\textbf{\ur{کل} \ur{مجھے} \ur{جانا} \ur{پڑا۔}}
\end{quote}

\begin{quote}
\emph{kal mujhe jānā paṛā.} — \textquotedblleft{}Yesterday I had to go.\textquotedblright{}
\end{quote}

\textbf{\ur{چاہیے}} is advice; \textbf{\ur{پڑنا}} is no choice.

\subsection*{Your turn}
Say these aloud:

\begin{itemize}
\raggedright
  \item \textquotedblleft{}I'll have to go\textquotedblright{} — \emph{mujhe jānā paṛe gā}
  \item \textquotedblleft{}I'll have to work\textquotedblright{} — \emph{mujhe kām karnā paṛe gā}
  \item \textquotedblleft{}I'll have to rest\textquotedblright{} — \emph{mujhe ārām karnā paṛe gā}
\end{itemize}

\begin{tcolorbox}[breakable,colback=teal!4,colframe=teal!35!black,title={Before you move on}]
\textquotedblleft{}You should go\textquotedblright{}? (\textbf{\ur{آپ} \ur{کو} \ur{جانا} \ur{چاہیے}}.) \textquotedblleft{}Yesterday I had to go\textquotedblright{}? (\textbf{\ur{کل} \ur{مجھے} \ur{جانا} \ur{پڑا}}.) Which one leaves no choice? (\textbf{\ur{پڑنا}}.)
\end{tcolorbox}
\section[\texorpdfstring{\ur{چنانچہ} (chunāṅche)}{چنانچہ (chunāṅche)}]{\ur{چنانچہ} (chunāṅche) — therefore}

\label{lesson:UR-C269-chunanche}



\begin{quote}
Before the new one: say \textquotedblleft{}so, that is why\textquotedblright{}. (\emph{is liye}.) Then \textquotedblleft{}I'll have to go\textquotedblright{}. (\emph{mujhe jānā paṛe gā}.)
\end{quote}

\subsection*{You'll want to know: \ur{چنانچہ}}
\textbf{\ur{چنانچہ}} — \emph{chunāṅche} — \textquotedblleft{}therefore\textquotedblright{}.

\begin{quote}
\textbf{\ur{بارش} \ur{ہو} \ur{رہی} \ur{تھی،}}
\end{quote}

\begin{quote}
\emph{bārish ho rahī thī,}
\end{quote}

\begin{quote}
\textbf{\ur{چنانچہ} \ur{ہم} \ur{گھر} \ur{پر} \ur{رہے۔}}
\end{quote}

\begin{quote}
\emph{chunāṅche ham ghar par rahe.} — \textquotedblleft{}It was raining; therefore we stayed at home.\textquotedblright{}
\end{quote}

It is formal: you will meet it in writing, in the news and in speeches. In talk, \textbf{\ur{اس} \ur{لیے}} (\emph{is liye}) says the same.

\begin{cousinweb}
\textbf{\ur{چنانچہ}} (\emph{chunāṅche}) is usually traced to Persian \emph{chunānche}, \textquotedblleft{}as, accordingly\textquotedblright{}, built on \emph{chunān}, \textquotedblleft{}so, thus\textquotedblright{}.
\end{cousinweb}

\subsection*{Your turn}
\begin{itemize}
\raggedright
  \item \emph{Say it:} \emph{chunāṅche}, then its everyday partner \emph{is liye}
  \item \emph{Read it:} \textbf{\ur{چنانچہ} \ur{ہم} \ur{گھر} \ur{پر} \ur{رہے}} aloud, as if from a newspaper
\end{itemize}

\begin{tcolorbox}[breakable,colback=teal!4,colframe=teal!35!black,title={Before you move on}]
What does \textbf{\ur{چنانچہ}} mean? (\textquotedblleft{}Therefore\textquotedblright{}.) Where do you meet it? (In writing and speeches.) Its everyday partner? (\textbf{\ur{اس} \ur{لیے}}.)
\end{tcolorbox}
\section[\texorpdfstring{\ur{اس} \ur{لیے} (is liye)}{اس لیے (is liye)}]{\ur{اس} \ur{لیے} (is liye) — so}

\label{lesson:UR-C269-is-liye}



\begin{quote}
Say \textquotedblleft{}because I was ill\textquotedblright{}. (\emph{kyūṅki maiṅ bīmār thā}.) That gives the reason. This lesson gives what follows from it.
\end{quote}

\subsection*{You'll want to know: \ur{اس} \ur{لیے}}
\begin{quote}
\textbf{\ur{اس} \ur{لیے} \ur{میں} \ur{نہیں} \ur{آیا۔}}
\end{quote}

\begin{quote}
\emph{is liye maiṅ nahīṅ āyā.} — \textquotedblleft{}So I didn't come.\textquotedblright{}
\end{quote}

\textbf{\ur{اس} \ur{لیے}} (\emph{is liye}), \textquotedblleft{}for this\textquotedblright{}, opens the result. The reason comes first, then the result:

\begin{quote}
\textbf{\ur{دیر} \ur{ہو} \ur{رہی} \ur{تھی،}}
\end{quote}

\begin{quote}
\emph{der ho rahī thī,}
\end{quote}

\begin{quote}
\textbf{\ur{اس} \ur{لیے} \ur{ہم} \ur{نے} \ur{ٹیکسی} \ur{لی۔}}
\end{quote}

\begin{quote}
\emph{is liye ham ne ṭaiksī lī.} — \textquotedblleft{}It was getting late, so we took a taxi.\textquotedblright{}
\end{quote}

\textbf{\ur{ٹیکسی}} (\emph{ṭaiksī}) is feminine: \textbf{\ur{لی}} (\emph{lī}).

\subsection*{You'll want to know: \ur{تو}, \textquotedblleft{}then\textquotedblright{}}
\begin{quote}
\textbf{\ur{باہر} \ur{بارش} \ur{ہے۔}}
\end{quote}

\begin{quote}
\emph{bāhar bārish hai.} — \textquotedblleft{}It's raining outside.\textquotedblright{}
\end{quote}

\begin{quote}
\textbf{\ur{تو} \ur{میں} \ur{گھر} \ur{پر} \ur{رہوں} \ur{گا۔}}
\end{quote}

\begin{quote}
\emph{to maiṅ ghar par rahūṅ gā.} — \textquotedblleft{}Then I'll stay at home.\textquotedblright{}
\end{quote}

\textbf{\ur{تو}} (\emph{to}) opens an answer that follows from what someone said: \textquotedblleft{}so, then\textquotedblright{}. In writing, \textbf{\ur{چنانچہ}} does the work of \textbf{\ur{اس} \ur{لیے}}.

\subsection*{Your turn}
Say these aloud:

\begin{itemize}
\raggedright
  \item \textquotedblleft{}I was ill, so I stayed at home\textquotedblright{} — \emph{maiṅ bīmār thā, is liye maiṅ ghar par rahā}
  \item \textquotedblleft{}It was raining, so I stayed at home\textquotedblright{} — \emph{bārish ho rahī thī, is liye maiṅ ghar par rahā}
  \item \textquotedblleft{}It was getting late, so I stayed at home\textquotedblright{} — \emph{der ho rahī thī, is liye maiṅ ghar par rahā}
\end{itemize}

\begin{tcolorbox}[breakable,colback=teal!4,colframe=teal!35!black,title={Before you move on}]
Which comes first, the reason or \textbf{\ur{اس} \ur{لیے}}? (The reason.) \textquotedblleft{}Then I'll stay at home\textquotedblright{}? (\textbf{\ur{تو} \ur{میں} \ur{گھر} \ur{پر} \ur{رہوں} \ur{گا}}.) And the formal word? (\textbf{\ur{چنانچہ}}.)
\end{tcolorbox}
\section[(review)]{Review — should, must, and so}

\label{lesson:UR-R269-advice-and-reasons}



\begin{quote}
Say \textquotedblleft{}I think it will rain\textquotedblright{}. (\emph{mujhe lagtā hai ki bārish ho gī}.) Then \textquotedblleft{}because of the rain\textquotedblright{}. (\emph{bārish kī vajah se}.)
\end{quote}

\subsection*{Your turn}
Say these aloud:

\begin{itemize}
\raggedright
  \item You should rest. You shouldn't go.
  \item I have to work every day.
  \item Yesterday I had to go to the office.
  \item I was ill, so I didn't come.
  \item It's raining? Then I'll stay at home.
\end{itemize}

\begin{tcolorbox}[breakable,colback=teal!4,colframe=teal!35!black,title={Before you move on}]
\textquotedblleft{}You should drink water\textquotedblright{}? (\textbf{\ur{آپ} \ur{کو} \ur{پانی} \ur{پینا} \ur{چاہیے}}.) Now the other way round: which word says there is no choice, \textbf{\ur{چاہیے}} or \textbf{\ur{پڑنا}}? (\textbf{\ur{پڑنا}}.) The formal word for \textbf{\ur{اس} \ur{لیے}}? (\textbf{\ur{چنانچہ}}.)
\end{tcolorbox}
